% =========================================================
% PART II -- DESCRIPTIVE STATISTICS
% =========================================================

\section{Descriptive Statistics}

% ---------------------------------------------------------
% 7. MEASURES OF CENTRAL TENDENCY
% ---------------------------------------------------------

\begin{frame}{7. Measures of Central Tendency}

\textbf{Definition}

Measures of central tendency describe the \textbf{central or typical
value} of a dataset.

\vspace{0.4cm}

The three commonly used measures are:

\begin{itemize}
    \item \textbf{Mean} -- arithmetic average of the observations.
    \item \textbf{Median} -- middle value after arranging the data.
    \item \textbf{Mode} -- most frequently occurring value.
\end{itemize}

\vspace{0.4cm}

\textbf{Purpose}

They provide a single value that represents the center of a dataset.

\vspace{0.3cm}

\textbf{Application:}

Used to summarize datasets in statistics, data analysis,
and machine learning.

\end{frame}


% ---------------------------------------------------------
% 7.1 MEAN -- DEFINITION AND FORMULA
% ---------------------------------------------------------

\begin{frame}{7.1 Mean}

\textbf{Definition}

The arithmetic mean is obtained by dividing the sum of all
observations by the number of observations.

\textbf{Formula}

For observations $x_1,x_2,\ldots,x_n$,
\[
\boxed{
\bar{x} = \frac{1}{n}\sum_{i=1}^{n}x_i
}
\]

where $n$ is the number of observations.

\textbf{Interpretation}

The mean represents the value obtained when the total of all
observations is distributed equally among the $n$ observations.

\textbf{Note:} The mean uses every observation in the dataset and
is therefore sensitive to extreme values.

\end{frame}


% ---------------------------------------------------------
% 7.1 MEAN -- EXAMPLE
% ---------------------------------------------------------

\begin{frame}{7.1 Mean -- Example}

\textbf{Example}

Consider the dataset:
\[
X=\{10,20,30,40,50\}
\]

There are $n=5$ observations.

\textbf{Calculation}
\[
\bar{x}
=
\frac{10+20+30+40+50}{5}
\]
\[
\boxed{\bar{x}=30}
\]

\textbf{Interpretation}

The average value of the dataset is $30$.

\textbf{Application}

The mean is widely used to summarize numerical features such as
marks, income, temperature, and sensor measurements.

\end{frame}


% ---------------------------------------------------------
% 7.2 MEDIAN -- DEFINITION AND FORMULA
% ---------------------------------------------------------

\begin{frame}{7.2 Median}

\textbf{Definition}

The median is the \textbf{middle value} of an ordered dataset.

The observations must first be arranged in ascending or
descending order.

\textbf{Formula}

For an ordered dataset with $n$ observations:
\[
\boxed{
\text{Median} =
\begin{cases}
x_{\frac{n+1}{2}}, & n \text{ is odd}\\[8pt]
\dfrac{x_{\frac{n}{2}}+x_{\frac{n}{2}+1}}{2},
& n \text{ is even}
\end{cases}
}
\]

\textbf{Key idea}

\begin{itemize}
    \item If $n$ is odd, the median is the single middle observation.
    \item If $n$ is even, the median is the average of the two middle observations.
\end{itemize}

\end{frame}


% ---------------------------------------------------------
% 7.2 MEDIAN -- EXAMPLES
% ---------------------------------------------------------

\begin{frame}{7.2 Median -- Examples}

\textbf{Example 1 -- Odd number of observations}

Consider:
\[
X=\{10,20,30,40,50\}
\]

Here, $n=5$ is odd.
\[
\text{Median}
=
x_{\frac{5+1}{2}}
=
x_3
=
\boxed{30}
\]

\vspace{0.1cm}

\textbf{Example 2 -- Even number of observations}

Consider:

\[
X=\{10,20,30,40\}
\]

Here, $n=4$ is even.

\[
\text{Median}
=
\frac{x_2+x_3}{2}
=
\frac{20+30}{2}
=
\boxed{25}
\]

\end{frame}


% ---------------------------------------------------------
% 7.3 MEDIAN -- OUTLIER EFFECT
% ---------------------------------------------------------

\begin{frame}{7.3 Median and Outliers}

\textbf{Example}

Consider:
\[
X=\{10,20,30,40,1000\}
\]

The mean is
\[
\bar{x}
=
\frac{10+20+30+40+1000}{5}
=
220
\]

while the median is
\[
\text{Median}=30
\]

\vspace{0.1cm}
\textbf{Interpretation}

The extreme value $1000$ greatly increases the mean, while the
median remains close to the center of the majority of observations.

\textbf{Key Point}

\begin{itemize}
    \item Mean is sensitive to extreme values.
    \item Median is more resistant to outliers.
\end{itemize}

\end{frame}


% ---------------------------------------------------------
% 7.4 MODE
% ---------------------------------------------------------

\begin{frame}{7.4 Mode}

\textbf{Definition}

The mode is the value that occurs with the \textbf{highest frequency}
in a dataset.

\textbf{Example}

Consider:
\[
X=\{2,4,4,5,6,4,7,8\}
\]

Therefore,
\[
\boxed{\text{Mode}=4}
\]
\textbf{Types}
\begin{itemize}
    \item \textbf{Unimodal} -- one mode.
    \item \textbf{Bimodal} -- two modes.
    \item \textbf{Multimodal} -- more than two modes.
\end{itemize}

\textbf{Application:}

Useful for identifying the most common category or value,
especially for categorical data.

\end{frame}


% ---------------------------------------------------------
% 7.5 COMPARISON
% ---------------------------------------------------------

\begin{frame}{7.5 Mean, Median and Mode}

\small

\begin{table}
\centering

\begin{tabular}{>{\bfseries}l l l}
\toprule
Measure & Main Idea & Sensitivity to Outliers \\
\midrule
Mean & Arithmetic average & High \\
Median & Middle value & Low \\
Mode & Most frequent value & Generally low \\
\bottomrule
\end{tabular}

\end{table}

\textbf{When to use?}

\begin{itemize}
    \item \textbf{Mean:} Suitable when the data is approximately
    symmetric and does not contain strong outliers.
    
    \item \textbf{Median:} Suitable for skewed data or data
    containing extreme values.
    
    \item \textbf{Mode:} Useful when the most frequent value or
    category is important.
\end{itemize}

\textbf{Application:}

Choosing an appropriate measure of central tendency helps
summarize real-world datasets effectively.

\end{frame}